\documentclass[a4paper, article, english, oneside, unknownkeysallowed, 12pt]{memoir} 
\usepackage{preamble}


\usepackage[
  backend=biber,
  style=authoryear-ibid,
  maxcitenames=2,  % show up to 2 authors in citations
  mincitenames=1,  % use "et al." if more than 2
  maxbibnames=10,  % show up to 10 names in bibliography
  minbibnames=1,
    uniquelist=false   
]{biblatex}


\addbibresource{causal_effect_schools.bib}

% Name and Title
\title{Causal effects of schools - March update}
\author{ Bruno E. Aravena-Magüida}
\date{March 2026}

% Header (top) and footer (bottom)
\pagestyle{fancy}
\fancyhead[R]{\emph{3rd-year micro group}}
 \fancyhead[L]{BEAM} % keep section names
\fancyfoot[R]{}
\fancyfoot[L]{}


\begin{document}

% Prints title block from macros given above
\maketitle

% Prints a table of contents
\setcounter{tocdepth}{2}
\renewcommand\contentsname{}
%\tableofcontents* 
%\clearpage

\section{Introduction}

\section{Conceptual Framework} \label{framework}

In this section I describe how I am currently thinking of my analysis.

\subsection{Potential Outcomes in Achievement}

For this section I borrow from \textcite{angrist_credible_2024} and \textcite{abdulkadiroglu_research_2017}

Suppose $N$ students and $J$ schools. 
Following the value-added literature on school effectiveness,  I assume a constant-effects model. 
This is equivalent to writing the potential outcome $Y_{ij}$ of a student $i$ in a school $j$ linearly separable as: 
\begin{equation}
  Y_{ij} = \nu_j + \epsilon_i
  \label{linear-separable}
\end{equation}

Hence, the causal effect of a school $j$ relative to a school $k$ would be given by $\nu_j - \nu_k$.


By adding a constant, we can renormalize this model relative to the average achievement ($\beta_0 = \frac{1}{J}\sum_{j=1}^J \nu_j$ )
As a result, we can rewrite the model as: 

\begin{equation}
 Y_i = \beta_0 + \sum_{j=1}^J \beta_j D_{ij} + \epsilon_i 
  \label{base}
\end{equation}

where now $\beta_j = \nu_j - \beta_0$. 

Unless school assignment is completely random, the evident threat to running this regression is the fact that $\cov(\epsilon_i, D_{ij})$ might be different than zero. 
This selection bias would refrain our estimates from having a causal interpretation. 

\subsubsection{Centralized Assignment Mechanisms}
A common strategy used in research to overcome this issue is the use of individual school lotteries, and more recently, centralized school assignments. 
In centralized assignment mechanisms, applying students submit rank-ordered preferences over schools ($\succ_i$) at which they might also have different priority levels\footnote{e.g. in the Chilean case being the child of a member of the school staff gives a student priority in that school.}.
These priorities operate as the school's preferences for students. An algorithm in place (most commonly student-proposing Deferred Acceptance) takes these preferences and outputs a school assignment ($Z_{ij}$) with one offer per student.
In order to break ties of students with the same preference and priority level, these systems randomize. 

If we refer to a student's priority at each school as $\rho_i = (\rho_{i1}, ...,\rho_{iJ})'$, we can define a student's type as the combination of both $\theta_i \equiv (\succ_i, \rho_i)$.
Let the probabilities of admission at a school be given by $p_i = (p_{i1}, ..., p_{ij})'$, with $p_{ij} = Pr(Z_{ij} = 1)$.

The mechanism described above implies that students with the same preferences and priority vector have the same probability of being assgined to the schools in their portfolio, which is called by \textcite{abdulkadiroglu_research_2017} the `equal treatment of equals' (ETE) property. 
Formally ETE states $\theta_i = \theta_j \implies p_i = p_j$. 

They also prove that under ETE that for any observed characteristic fixed under re-randomization $W_i$ (i.e. pre-assignment or not affected by assignment) we have that:
\begin{equation}
 Pr[Z_{ij} = 1|\, W_i = w, \theta_i = \theta] = Pr[Z_{ij} = 1 | \, \theta_i = \theta]  \coloneqq p_j(\theta) 
\end{equation}

Hence, for any such variable $W_i$ we can write 
\begin{equation}
 W_i \ind Z_i \, | \, \theta_i 
 \label{w ind}  
\end{equation}

In particular, and as a consequence of equation \ref{linear-separable}, where we assumed ability $\epsilon_i$ to be fixed under re-randomization, partly due to the absence of match effects, we can write
\begin{equation}
  \epsilon_i \ind Z_i \, | \, \theta_i   
\end{equation}

This is the core assumption of the design. 
Under this assumption and perfect compliance ($Z_i = D_i$), we could estimate equation \ref{base} by saturating on indicator for each student type $\theta_i$. 

However, that is often not possible as these strata might be too granular to have enough variation within each of them to estimate the model. 
As described by \textcite{angrist_credible_2024}: "In practice (...) we see nearly as many preference and priority combinations as there are students in a match."  
A solution to this issue is to use the probablities of admissions as propensity scores to pool over student types, as proposed by \textcite{abdulkadiroglu_research_2017}. This idea is just an application of the propensity score matching \parencite{rosenbaum_central_1983}, and it yields:  
\begin{equation}
  \epsilon_i \ind Z_i \, | \, p_i 
  \label{CRA}
\end{equation}

Note that by eq. \ref{w ind} we can also consider other control variables (as long as they are fixed under re-randomization). Thus for any pre-treatment covariate $X_i$ it must be the case that 
($\epsilon_i, X_i) \ind Z_i | \theta_i$, so we can rewrite Eq \ref{CRA} as $  \epsilon_i \ind Z_i \, | \, (p_i, X_i) $ to allow for such controls.

\subsubsection{RC VAM}

The traditional IV framework would allow us to estimate a causal effect with an assumption such as \ref{CRA} for any school where we have a non-degenerate lottery (an instrument -lottery based offer- for each endogenous -attendance- variable).
However, that would not allow to get an estimate for schools without lotteries. An hybrid approach proposed by (paper) suggests integrating the variation given by lotteries within a simpler value-added framework.

To do this, the authors make an addtional assumption, which is to have compliance with the offers to be as good as random (under the same conditioning set.)
They argue this is simlar to Dale and Krueger's strategy (quote req), where the sorting occurs when application choices are made and deviations from such plans must be unrelated to the potential outcomes.

This final assumption yields a traidional conditional independence assumption, from where estimation is straight-forward. 

\begin{equation}
  \epsilon_i \ind D_i \, | \, (p_i, X_i) 
  \label{CIA}
\end{equation}

In particular, if covariates used are lagged (pre-treatment) outcomes we end up with a Risk Controlled Value Added Model (RC VAM) as constructed in (centralized VAM paper). 

%\begin{equation}
%  Y_i = \alpha_0 + \sum_j \alpha_j D_{ij} + X_i' \Gamma + g(p_i) + \eta_i 
%\end{equation}


\subsection{Estimation in specific outcomes}

In this subsection I present the main regressions I intend to run, some challenges and how I intend to adress them. 
Let $Y^{MATH}_i$ represent the outcome of a student in the standarized exam for mathematics and $Y^{STEM}_i$ an indicator variable for whether the student enrolls in a STEM program. 
Let $X_i$ be a vector of characteristics not changed by treatment (demographics), including also pre-treatment math achievement $Y^{MATH}_{pre,i}$.

Following equation \ref{CIA} I write:

\begin{equation}
  Y^{MATH}_i = b_0 + \sum_j b_j D_{ij} +  X_i' B +   g(p_i) + \epsilon_i 
  \label{Math}
\end{equation}


\begin{equation}
  Y^{STEM}_i = c_0 + \sum_j c_j D_{ij} + X_i' C + g(p_i) + \eta_i 
  \label{STEM_prev}
\end{equation}

The only difference between these two regressions is the fact that I propose to include pre-treatment achievement on the first regression and exclude it from the second. 
The incllusion of lagged outcome for math achievement makes \ref{Math} a VAM regression and it helps to increase precision. 
I do not have a similar lagged outcome for STEM enrollment. 
Remember that, under the CIA assumption above, these regressions could be run without any of the controls and should still hold causal interpretation: the CIA assumption only needs the $p_i$ to be included.


Consider instead in the true model for STEM enrollment we have the student outcome in the standarized math exam $Y^{MATH}_i$ entering linearly, and that everything else is specified as in equation \ref{STEM_prev}.
In such case, the true model would be given by: 
\begin{equation}
  Y^{STEM}_i = \gamma_0 + \sum_j \gamma_j D_{ij} +\gamma_M Y^{MATH}_i + X_i' \Gamma + g(p_i) + \eta_i 
  \label{STEM_true}
\end{equation}

This idea is based on the observed relationship (figure), which likely captures potentially both issues of elegibility into more mathematically-demanding programs, as well as sorting on related abilities. 


Then, if we run the regression \ref{STEM_prev} we would get that: 
\begin{align*}
\cov(Y^{STEM}_i; D_{ik} | X_i, p_i) = \cov(\gamma_0 + \sum_j \gamma_j D_{ij} +\gamma_M Y^{MATH}_i + X_i' \Gamma + g(p_i) + \eta_i; D_{ik} | X_i, p_i )   \\
= \cov(\gamma_0 + \sum_j \gamma_j D_{ij} +\gamma_M \big[  b_0 + \sum_j b_j D_{ij}  + X_i' B +   g(p_i) + \epsilon_i \big] + X_i' \Gamma + g(p_i) + \eta_i; D_{ik} | X_i, p_i ) \\
= \cov(\gamma_k D_{ik} +\gamma_M \big[  b_k D_{ik}  + \epsilon_i   \big] + \eta_i , D_ij | X_i, p_i) \\
= \cov(\gamma_k D_{ik} +\gamma_M \big[  b_k D_{ik}    \big] + ; D_{ik} | X_i, p_i) \\
= \big( \gamma_k + \gamma_M (b_k) \big) \big( \var(D_{ik} | X_i, p_i) \big) \\
\neq \gamma_k \big( \var(D_{ik} | X_i, p_i) \big) 
\end{align*}

Where the first equality follows from replacing equation \ref{Math} in \ref{STEM_true}. The second one from dropping constants and functions of the conditioning set. 
The third one from applying equation \ref{CIA} to the individual's fixed mathematical ability ($\epsilon_i$) and interest for STEM ($\eta_i$), and the final one from covariance/variance identities. 

Note that unless either $\gamma_M = 0$ (math scores not affecting likelihood of student applying to STEM) or $(b_k) = 0$ (school not affecting math scores), the outcome of the short regression is biased. 

A solution to this issue could be to include instead a forecast of $\hat{Y^{MATH}_i}$ constructed from pre-treatment variables such as $Y^{MATH}_{pre,i}$, potentially also including other elements of $X_i$.

\clearpage
\section{Structural model of major choice}

In order to separate these two channels, I intend to estimate a model of demand for fields following (Do Parents Value School Effectiveness?† 2020) and more closely (Campos et al. 2026). 
I folow more closely the work by (Campos et al. 2026)  as I decide to collapse the demand into broad fields of occupation as opposed to particular programs.

Let $u_{if}$ be the indirect utility associated with field $f \in \mathcal{F}$. Let $g(i)$ index the covariate-cell group of $i$. 
I specify the following functional form: 

\begin{equation}
  u_{if} = \delta_{g(i)f} - \kappa_{g(i)} d_{if} + f_{g(i)f}(Y^{MATH}, Y^{VERBAL}) + \lambda_{fs} + \eta_{if}  
\end{equation}

where $\delta_{g(i)f}$ represents the mean utility of field $j$ for applicants of covariate-cell $g(i)$, $d_{if}$ is given by the distance between student $i$ and the nearest program corresponding to the field $f$, and $\kappa_{g(i)}$ is a cell-specific distance cost. 
$f_{g(i)f}(Y^{MATH}, Y^{VERBAL})$ gives a polynomial of degree 2 for the math and verbal exams, for each covariate-cell and field. $\lambda_{fs}$ is a school-field preference shifter. Finally $\eta_{if}$ represents the unobserved preference shock, assumed to follow an EVT1 distribution conditional on $(X_i', d_i')$. 



The idea underlying this approach is to use the submitted college application ranked-choice lists as revealing their preferences given that in this setting truthful reporting is a weakly dominant strategy given that these also operate under a student-proposing DA system.
If we refer to a student's rank-order list by $R_i = (R_{i1}, ..., R_{il(i)})'$ where $l(i)$ is the number of different fields ranked by the student, 
then we have that the students' top-ranked field satisfies we have that:
\[R_{1i} = \arg \max_{f \in \mathcal{F}} U_{if} \] 

with the remaining options satisfying 
\[R_{ki} = \arg \max_{f \in \mathcal{F}  \textbackslash{} (R_{i1}, ..., R_{k-1}) } U_{if} \] 

The EVT1 assumption implies the following conditional likelihood of the rank list $R_i$: 
\begin{equation}
  \mathcal{L}(R_i | \, X_i, d_i) = \prod_{k=1}^{l(i)} \frac{\delta_{g(i)f} - \kappa_{g(i)} d_{if} + f_{g(i)f}(Y^{MATH}, Y^{VERBAL}) + \lambda_{fs}}{\sum_{r \in R_i \textbackslash{} R_{i1}, ..., R_{k-1}} \delta_{g(i)f} - \kappa_{g(i)} d_{if} + f_{g(i)f}(Y^{MATH}, Y^{VERBAL}) + \lambda_{fs} }
\end{equation}


For the covariate-cells, I consider each combination of macro-region (3 values), gender (2 values) and family income (above, below median: 2 values) for a total of 12 covariate-cells. For each of them, I need to compute 5 parameters for each field -mean utility and 2 polynomials of degree 2- plus an extra term for the sensitivity to distance. 
That would yield a total of 12 x (5(8) + 1) = 12(41) = 492 parameters. On top of that estimating school-preference shifter with 200 schools, would require other (200)(8) = 1600 parameters, for a total of 2092 parameters.
Limiting myself to study only 100 schools would bring down the number of parameters to 1292.

An important caveat is that this approach does not rely at all in the use of lotteries. As a result, as long as preferences are relatively estable over time, I could use data from cohorts that arrived to high school before the introduction of the lotteries. Potentially, this is also an interesting space to work on. 

\clearpage
\section{Computing $p_{ij}$} 

Evidently, as discussed in section \ref{framework}, the computation of the $p_{ij}$ is crucial to leveraging the exogenous variation provided by the centralized assignment mechanism.


\subsection{Reproducing school offers in the Chilean system}

I have made significant progress in being able to fully replicate the assignment delivered by the mechanism. On one hand this is possible by the fact that the responsible agency includes the used lottery number when making the data available. 
On the other hand, this process has become more complicated given the number of features of the system that escape a traditional DA mechanism. Documentation for these is far from complete.

The Chilean system has many features that are not captured by the textbook DA mechanism which features ranked preferences and priorities $(\succ_i, \rho_i)$.
There are reserved spots for low-income, high-achievement and special-needs students. Despite documentation not being clear about this,
these quotas do roll-over to regular spots if they are not taken. Additionally, students have guaranteed admission in their original school if they do not get assigned anywhere else, even if their original school did not list any available spots.
The system also allows siblings to apply in block, such that if the older one gets accepted by a school, the submitted list of younger siblings gets updated and prioritizes such school. Note this is different than the priority tier for students who have a sibiling at the school before the process starts. 
Some of these rules also change over a transition period on the first few years of the implementation, which is where I sample from given that this allows me to have enough cohorts that have spent at least 4 years after their centralized assignment.

I can currently replicate 99.47 \% of the assignments for almost 80k students applying to grade 9 in 2018.
To get a sense of the importance of getting these additional features right, note that using only the the individual's types and a random seed replicates only up to 68 \% of the assignments.
Including the actual lottery numbers in this case, slightly decreases the fraction of matches.

Evident next steps in this part of the project include trying to figure out the remaining fraction, as well as ensure that I can achieve similar replication rates for later admission cycles, where some of the transition rules phase out.


\subsection{Computing $p_{ij}$}

For a traditional DA system (MCM)



\section{Conclusion}

Bye bye 


\printbibliography


\end{document}